% from ega4-2-fr.tex
\oldpage[IV-2]{208}

Il est immédiat que cet homomorphisme est bijectif lorsque $M$ est libre de
type fini; comme on a une suite exacte $L'\to L\to M\to0$, où $L$ et $L'$ sont
des $A$-modules libres de type fini, on en déduit un diagramme commutatif
\[
\xymatrix{
L'\widehat\otimes_kB \ar[r] \ar[d] &
L\widehat\otimes_kB \ar[r] \ar[d] &
M\widehat\otimes_kB \ar[r] \ar[d] & 0 \\
L'\otimes_A(A\widehat\otimes_kB) \ar[r] &
L\otimes_A(A\widehat\otimes_kB) \ar[r] &
M\otimes_A(A\widehat\otimes_kB) \ar[r] & 0
}
\]
où les lignes sont exactes (la première en vertu de la définition du produit
tensoriel complété et de (0$_{\mathrm{III}}$, 13.2.2)); les deux premières
flèches verticales étant des bijections, il en est de même de la troisième.
\end{proof}

Le fait que $A\widehat\otimes_kB$ s'identifie à un sous-anneau de
$A'\otimes_A(A\widehat\otimes_kB)$ résulte alors de ce que $A$ est un
sous-anneau de $A'$ et de ce que $A\widehat\otimes_kB$ est un $A$-module plat
(7.5.5). On peut ainsi supposer en outre que $A$ et $B$ sont intégralement
clos; comme en outre $k$ est supposé parfait, $\operatorname{Spec}(A)$ et
$\operatorname{Spec}(B)$ sont géométriquement normaux sur $k$ en vertu de
(6.7.7, b)); on peut alors appliquer (7.5.6) en prenant pour
$\boldsymbol R$ la propriété (vi) de (7.5.3).

II) Supposons $k$ algébriquement clos, $A$ et $B$ intègres. Le raisonnement de
I) ramène encore au cas où $A$ et $B$ sont intégralement clos; il résulte alors
de (7.5.6) que $\operatorname{Spec}(A\widehat\otimes_kB)$ est normal, donc
$A\widehat\otimes_kB$ est composé direct d'anneaux locaux complets
intégralement clos, et tout revient à voir que $A\widehat\otimes_kB$ est un
anneau local; il suffit pour cela de vérifier que si $\mathfrak m$ et
$\mathfrak n$ sont les idéaux maximaux de $A$ et $B$,
$(A/\mathfrak m)\otimes_k(B/\mathfrak n)$ est un anneau local (voir
démonstration de (0, 19.7.1.2)); mais comme $A/\mathfrak m$ et
$B/\mathfrak n$ sont des extensions finies de $k$, elles sont identiques à
$k$, d'où la conclusion.
\end{proof}

\begin{remark}[7.5.8]
\label{IV.2.s7.5.8-fr}
Il serait intéressant de déterminer si, dans (7.5.7), on peut remplacer
l'hypothèse que $k$ est parfait ou algébriquement clos par l'hypothèse que
$\operatorname{Spec}(A)$ ou $\operatorname{Spec}(B)$ est géométriquement
intègre sur $k$, ou tout au moins par l'hypothèse que $A$ (par exemple) soit
intègre et contienne un sous-anneau $A_0$ isomorphe à un anneau de séries
formelles $k[[T_1,\ldots,T_n]]$, et que le corps des fractions $K$ de $A$ soit
séparable sur le corps des fractions $K_0$ de $A_0$ (on peut montrer que cette
condition entraîne que $A$ est géométriquement intègre sur $k$, et que les
deux propriétés sont équivalentes lorsque $[k:k^p]$ est fini ($p$ étant
l'exposant caractéristique de $k$)). De même il serait désirable de développer
des variantes de (7.5.6) et (7.5.7) où on affaiblirait l'hypothèse de finitude
des corps résiduels de $A$ et $B$, en supposant par exemple qu'un seul d'entre
eux est extension finie de $k$, l'autre étant quelconque.
\end{remark}

\subsection{Applications : I. Anneaux japonais locaux}
\label{IV.2.ss7.6-fr}

\begin{proposition}[7.6.1]
\label{IV.2.s7.6.1-fr}
Soit $A$ un anneau local noethérien réduit dont les fibres formelles sont
géométriquement normales. Alors le complété $\widehat A$ est réduit, la
fermeture intégrale $A'$ de $A$ dans son anneau total des fractions est une
$A$-algèbre finie (donc un anneau noethérien semi-local) et son complété
$\widehat{A'}$ est isomorphe à la fermeture intégrale de $\widehat A$ dans son
anneau total des fractions.
\end{proposition}

\begin{proof}
Les fibres formelles de $A$ sont a fortiori géométriquement réduites, donc
l'hypothèse que $A$ est réduit entraîne qu'il en est de même de $\widehat A$,
en vertu de (7.3.17). Soient $\mathfrak p_i$ ($1\leq i\leq n$) les idéaux
premiers minimaux de $A$, et soit $B_i=A/\mathfrak p_i$; les fibres formelles
des anneaux locaux $B_i$ sont alors aussi géométriquement normales (7.3.15),
donc les $\widehat B_i$ sont aussi réduits et il résulte de (0, 23.1.7, (i))
que la clôture intégrale $B'_i$ de $B_i$ dans son corps des fractions est un
$B_i$-module de type fini, donc un $A$-module de type fini; comme $A'$ est
composé direct des $B'_i$ (II, 6.3.8), on voit que $A'$ est un $A$-module de
type fini. Soient $\mathfrak m_j$ ($1\leq j\leq r$) les idéaux maximaux de
l'anneau semi-local $A'$; on sait que le complété $\widehat{A'}$ de $A'$
s'identifie au composé direct des complétés
$\widehat{A'_{\mathfrak m_j}}$ des $A'_{\mathfrak m_j}$ (Bourbaki,

